\documentclass[10pt,a4paper]{amsart}
\usepackage{amsmath,amssymb,amsthm,graphicx}
% Retained minimal preamble for the Stacks Project chapter more-algebra.tex (More on Algebra).
% The chapter file pinned-original-source.tex carries only \input{preamble} on its first line, so no
% preamble text lives inside the chapter itself.  The Stacks Project uses one single shared preamble for
% every chapter; the snapshot stages it identically as sources/stacks-chow/preamble.tex and
% sources/stacks-etale-cohomology/preamble.tex.  The block below is copied mechanically, line for line,
% from sources/stacks-chow/preamble.tex lines 14-18 (the Xy-pic setup) and 154-217 (the native theorem
% declarations together with the mathematical macros).  Nothing was added, removed or redefined; the
% wrapper class and the amsmath/amssymb/amsthm/graphicx base packages are supplied by the assembler.
% Chapter: more-algebra.tex -- More on Algebra.
% For commutative diagrams we use Xy-pic
\usepackage[all]{xy}

% We use 2cell for 2-commutative diagrams.
\xyoption{2cell}

% Theorem environments.
%
\theoremstyle{plain}
\newtheorem{theorem}[subsection]{Theorem}
\newtheorem{proposition}[subsection]{Proposition}
\newtheorem{lemma}[subsection]{Lemma}

\theoremstyle{definition}
\newtheorem{definition}[subsection]{Definition}
\newtheorem{example}[subsection]{Example}
\newtheorem{exercise}[subsection]{Exercise}
\newtheorem{situation}[subsection]{Situation}

\theoremstyle{remark}
\newtheorem{remark}[subsection]{Remark}
\newtheorem{remarks}[subsection]{Remarks}

\numberwithin{equation}{subsection}

% Macros
%
\def\lim{\mathop{\mathrm{lim}}\nolimits}
\def\colim{\mathop{\mathrm{colim}}\nolimits}
\def\Spec{\mathop{\mathrm{Spec}}}
\def\Hom{\mathop{\mathrm{Hom}}\nolimits}
\def\Ext{\mathop{\mathrm{Ext}}\nolimits}
\def\SheafHom{\mathop{\mathcal{H}\!\mathit{om}}\nolimits}
\def\SheafExt{\mathop{\mathcal{E}\!\mathit{xt}}\nolimits}
\def\Sch{\mathit{Sch}}
\def\Mor{\mathop{\mathrm{Mor}}\nolimits}
\def\Ob{\mathop{\mathrm{Ob}}\nolimits}
\def\Sh{\mathop{\mathit{Sh}}\nolimits}
\def\NL{\mathop{N\!L}\nolimits}
\def\CH{\mathop{\mathrm{CH}}\nolimits}
\def\proetale{{pro\text{-}\acute{e}tale}}
\def\etale{{\acute{e}tale}}
\def\QCoh{\mathit{QCoh}}
\def\Ker{\mathop{\mathrm{Ker}}}
\def\Im{\mathop{\mathrm{Im}}}
\def\Coker{\mathop{\mathrm{Coker}}}
\def\Coim{\mathop{\mathrm{Coim}}}

% Boxtimes
%
\DeclareMathSymbol{\boxtimes}{\mathbin}{AMSa}{"02}

%
% Macros for moduli stacks/spaces
%
\def\QCohstack{\mathcal{QC}\!\mathit{oh}}
\def\Cohstack{\mathcal{C}\!\mathit{oh}}
\def\Spacesstack{\mathcal{S}\!\mathit{paces}}
\def\Quotfunctor{\mathrm{Quot}}
\def\Hilbfunctor{\mathrm{Hilb}}
\def\Curvesstack{\mathcal{C}\!\mathit{urves}}
\def\Polarizedstack{\mathcal{P}\!\mathit{olarized}}
\def\Complexesstack{\mathcal{C}\!\mathit{omplexes}}
% \Pic is the operator that assigns to X its picard group, usage \Pic(X)
% \Picardstack_{X/B} denotes the Picard stack of X over B
% \Picardfunctor_{X/B} denotes the Picard functor of X over B
\def\Pic{\mathop{\mathrm{Pic}}\nolimits}
\def\Picardstack{\mathcal{P}\!\mathit{ic}}
\def\Picardfunctor{\mathrm{Pic}}
\def\Deformationcategory{\mathcal{D}\!\mathit{ef}}

% --- appended 2026-10-12 (Chief): the project own definitions of the reference/slogan/history
% display environments, copied mechanically from the shared Stacks preamble lines 8-12
% (\usepackage{verbatim} plus the three \newenvironment lines, which the chapter preamble keeps
% verbatim from the project preamble). Nothing added, removed or redefined: these environments
% carry no mathematical content and the project itself typesets them as comments. Without these
% three lines a counted statement whose own text uses them fails with Environment slogan undefined.
\usepackage{verbatim}
\newenvironment{reference}{\comment}{\endcomment}
\newenvironment{slogan}{\comment}{\endcomment}
\newenvironment{history}{\comment}{\endcomment}
\usepackage{cleveref}
\crefname{originallemma}{original lemma}{original lemmas}
\Crefname{originallemma}{Original lemma}{Original lemmas}
\makeatletter
\newlabel{algebra-lemma-henselization-different}{{04GV}{source}{}{sourcepointer-1}{}}
\newlabel{algebra-lemma-henselization-different@cref}{{[originallemma][0][]04GV}{[1][0][]source}{}{}{}}
\makeatother
\crefname{originallemma}{original lemma}{original lemmas}
\Crefname{originallemma}{Original lemma}{Original lemmas}
\makeatletter
\newlabel{lemma-etale-extension-valuation-ring}{{0ASJ}{source}{}{sourcepointer-2}{}}
\newlabel{lemma-etale-extension-valuation-ring@cref}{{[originallemma][0][]0ASJ}{[1][0][]source}{}{}{}}
\makeatother
\crefname{originallemma}{original lemma}{original lemmas}
\Crefname{originallemma}{Original lemma}{Original lemmas}
\makeatletter
\newlabel{algebra-lemma-colimit-valuation-rings}{{0AS4}{source}{}{sourcepointer-3}{}}
\newlabel{algebra-lemma-colimit-valuation-rings@cref}{{[originallemma][0][]0AS4}{[1][0][]source}{}{}{}}
\makeatother
\crefname{originallemma}{original lemma}{original lemmas}
\Crefname{originallemma}{Original lemma}{Original lemmas}
\makeatletter
\newlabel{algebra-lemma-strict-henselization-different}{{04GW}{source}{}{sourcepointer-4}{}}
\newlabel{algebra-lemma-strict-henselization-different@cref}{{[originallemma][0][]04GW}{[1][0][]source}{}{}{}}
\makeatother
\title{Henselization and strict henselization of a valuation ring are weakly unramified extensions}
\author{The Stacks Project Authors}
\date{stacks-project:more-algebra.tex:local-snapshot-20261012; GNU FDL 1.2 or later}
\begin{document}
\maketitle
\setcounter{subsection}{1}
\noindent\textit{Editorial note.} All delivered mathematical text is the author's own, copied from the pinned Stacks Project chapter file more-algebra.tex; no statement or proof step was rewritten, reordered or removed. The counted statement is the labeled lemma environment at source lines 37534-37543 (canonical tag 0ASK) and the counted proof is the proof environment at source lines 37545-37566, which immediately follows it and is its own proof.\par\medskip
\noindent\textit{Editorial note.} The statement and the proof display the chain of inclusions and the map on value groups; both displays are copied exactly as printed, using the macros of the retained preamble.\par\medskip
\noindent\textit{Editorial note.} The four references are one to a result of this chapter, namely the statement about etale extensions of valuation rings, and three to results of other chapters on descriptions of henselization and on colimits of valuation rings; the kinds recorded are the type words appearing in the labels themselves. All four are navigation stubs with their roles recorded in the reasons.\par\medskip
\noindent\textit{Editorial note.} The preamble span runs to line 90, the current end of the retained preamble, which includes the appended environment definitions; the counted text does not use them, and the span is taken whole so that the file version is unambiguous.\par\medskip
\noindent\textit{Editorial note.} The source URL is the source-level Stacks Project URL, as required when more than one card is produced from one source; the per-claim tag is carried by claim\_key and by this note. The counted statement and proof cite no bibliography entry, so this card needs no bibliography.\par\medskip
\noindent\textit{Allowed standard prerequisites (editorial).}\begin{itemize}
\item The henselization of a local ring is the filtered colimit of the localizations of etale algebras at primes lying over the maximal ideal, and strict henselization is described by the analogous statement. invoked at proof lines 37545-37549 and again at proof lines 37562-37564. Produces the colimit presentation on which the whole computation rests.
\item The localization of an etale algebra over a valuation ring at a prime over the maximal ideal is again a valuation ring, and the map on value groups is bijective. invoked at proof lines 37549-37556. Gives the statement for each term of the colimit, including the bijectivity of the map on value groups.
\item A filtered colimit of valuation rings with compatible value groups is a valuation ring, and the value group of the colimit is computed from the colimit of the value groups. invoked at proof lines 37556-37561. Passes from the terms of the colimit to the henselization itself, and then to the strict henselization and the composite inclusion.
\end{itemize}\medskip
\ifdefined\hypertarget\hypertarget{sourcepointer-1}{}\fi
\noindent\textit{Source reference.} Original lemma 04GV: sources/stacks-chow/tags.txt, source lines 5846--5846. The reference is the description of the henselization as a colimit of etale localizations, used at proof lines 37545-37549; the computation performed on that colimit is inside the counted span.\par\medskip
\ifdefined\hypertarget\hypertarget{sourcepointer-2}{}\fi
\noindent\textit{Source reference.} Original lemma 0ASJ: sources/stacks-more-algebra/pinned-original-source.tex, source lines 37480--37480. The reference is applied at proof lines 37549-37556 to each term of the colimit to get a valuation ring with the same value group; it is the input to the colimit passage, which is inside the counted span.\par\medskip
\ifdefined\hypertarget\hypertarget{sourcepointer-3}{}\fi
\noindent\textit{Source reference.} Original lemma 0AS4: sources/stacks-chow/tags.txt, source lines 13567--13567. The reference passes the property from the terms to the filtered colimit, used at proof lines 37556-37561; the conclusion about value groups drawn there is inside the counted span.\par\medskip
\ifdefined\hypertarget\hypertarget{sourcepointer-4}{}\fi
\noindent\textit{Source reference.} Original lemma 04GW: sources/stacks-chow/tags.txt, source lines 5847--5847. The reference is the strict analogue of the colimit description, used at proof lines 37562-37564 so that the same argument applies to the strict henselization; the identification of the composite inclusion is inside the counted span.\par\medskip
\begin{lemma}
\label{lemma-henselization-valuation-ring}
Let $A$ be a valuation ring. Let $A^h$, resp.\ $A^{sh}$ be its
henselization, resp.\ strict henselization. Then
$$
A \subset A^h \subset A^{sh}
$$
are extensions of valuation rings which induce bijections on
value groups, i.e., which are weakly unramified.
\end{lemma}

\begin{proof}
Write $A^h = \colim (B_i)_{\mathfrak q_i}$ where $A \to B_i$
is \'etale and $\mathfrak q_i \subset B_i$ is a prime ideal
lying over $\mathfrak m_A$, see
Algebra, Lemma \ref{algebra-lemma-henselization-different}.
Then Lemma \ref{lemma-etale-extension-valuation-ring}
tells us that $(B_i)_{\mathfrak q_i}$
is a valuation ring and that the induced map
$$
(A \setminus \{0\})/A^* \longrightarrow
((B_i)_{\mathfrak q_i} \setminus \{0\}) / (B_i)_{\mathfrak q_i}^*
$$
is bijective. By Algebra, Lemma \ref{algebra-lemma-colimit-valuation-rings}
we conclude that $A^h$ is a valuation ring. It also follows that
$(A \setminus \{0\})/A^* \to (A^h \setminus \{0\})/(A^h)^*$
is bijective. This proves the lemma for the inclusion $A \subset A^h$.
To prove it for $A \subset A^{sh}$ we can use exactly the same argument
except we replace Algebra, Lemma \ref{algebra-lemma-henselization-different} by
Algebra, Lemma \ref{algebra-lemma-strict-henselization-different}.
Since $A^{sh} = (A^h)^{sh}$ we see that this also proves the
assertions of the lemma for the inclusion $A^h \subset A^{sh}$.
\end{proof}

\end{document}
